\section{Anlageurteil}
\label{sec:urteil-azn}

\subsection{Was belegt ist}

Die schwarz gesetzten Befunde dieses Teils in einem Absatz, ohne Wertung: Der
Umsatz stieg 2025 um \Pct{\AznUmsatzWachstum} auf \MioUSD{\AznUmsatzXXV}, der
Jahres\"uberschuss um \Pct{\AznErgebnisWachstum} auf
\MioUSD{\AznJahresueberschussXXV}. Die EBITDA-Marge betr\"agt
\Pct{\AznEbitdaMarge}, die Eigenkapitalrendite \Pct{\AznEigenkapitalrendite},
die Nettoverschuldung das \Faktor{\AznVerschuldungsgrad}-fache des EBITDA. Der
freie Mittelzufluss je Aktie erreichte \USDje{\AznFcfVoll} und damit den
h\"ochsten Wert der sieben erfassten Jahre. Die Dividende von
\USDje{\AznDividendeJeAktie} ist \Faktor{\AznDividendeDeckungCore}-fach durch
das bereinigte Ergebnis gedeckt; R\"uckk\"aufe finden nicht statt. Das
Unternehmen hat seine Umsatzprognose in vier belegbaren Jahren dreimal
getroffen und einmal \"ubertroffen. Das gr\"o{\ss}te Einzelmedikament mit
\Pct{\AznFarxigaAnteil} des Umsatzes hat im zweiten Quartal 2026 in den
Vereinigten Staaten den Patentschutz verloren. Der Kurs steht
\PctBetrag{\AznAbstandHoch} unter seinem Jahreshoch.

\subsection{Was das Urteil tr\"agt}

\bk{Drei Beobachtungen tragen das Urteil, und sie zeigen in
unterschiedliche Richtungen.}

\bk{\emph{Die Qualit\"at des Gesch\"afts ist nicht die Frage.} Eine
Eigenkapitalrendite von \Pct{\AznEigenkapitalrendite}, eine EBITDA-Marge von
\Pct{\AznEbitdaMarge}, eine Verschuldung unter dem Eineinhalbfachen des EBITDA
und eine Prognose, die in vier Jahren nie unterschritten wurde: Das ist ein
Unternehmen, an dem sich wenig auszusetzen findet. Wer diesen Bericht liest,
um zu erfahren, ob AstraZeneca gut gef\"uhrt ist, kann hier aufh\"oren.}

\bk{\emph{Die Frage ist der Preis.} Bei \USDje{\AznKurs} zahlt der K\"aufer
das \Faktor{\AznKursGewinnCore}-fache des bereinigten Gewinns und das
\Faktor{\AznKursBuchwert}-fache des Buchwerts. Der freie Mittelzufluss je
Aktie von \USDje{\AznFcfVoll} entspricht einer Rendite von
\Pct{\AznFcfRendite} auf den Kurs. Damit daraus \Pct{\Huerde} pro Jahr werden,
muss entweder der Mittelzufluss wachsen oder die Bewertung steigen; auf dem
heutigen Niveau stehenbleiben reicht nicht.}

\bk{\emph{Das Basisjahr ist das beste der Reihe.} Das ist die Beobachtung, die
am leichtesten \"ubersehen wird und am meisten wiegt. Alle
Renditerechnungen dieses Teils gehen von einem Niveau aus, das im hundertsten
Perzentil der eigenen sieben Jahre liegt. Ein Modell, das den Spitzenwert
konstant fortschreibt, ist bereits optimistisch, und es erreicht die H\"urde
trotzdem nicht. Die Gegenprobe mit dem Median (\USDje{\AznFcfMedian}) und mit
den letzten zw\"olf Monaten (\USDje{\AznFcfLtm}) f\"allt entsprechend
ung\"unstiger aus.}

\subsection{Nicht verf\"ugbare Angaben}

Vier Gr\"o{\ss}en fehlen, und keine davon wird gesch\"atzt.

Der Umsatz, der insgesamt an auslaufenden Patenten h\"angt, ist vom
Unternehmen nicht beziffert (Abschnitt~\ref{sec:loe-azn}). Das ist die
gr\"o{\ss}te L\"ucke dieses Teils, weil sie den Nenner der wichtigsten Frage
betrifft.

Die Zusammensetzung des Halbjahresergebnisses nach Therapiegebieten liegt in
gr\"oberer Form vor als die des Jahres; eine Fortschreibung einzelner
Produktreihen auf das Gesamtjahr w\"are eine Sch\"atzung und unterbleibt.

Die Prognose f\"ur das Gesch\"aftsjahr 2021 ist nicht belegbar
(Abschnitt~\ref{sec:prognose-azn}); die Prognosetreue umfasst deshalb vier
statt f\"unf Jahre.

Ein Wechselkurs zum Stichtag der Verg\"utungstabelle fehlt; die Angaben in
Pfund bleiben unumgerechnet.
